\begin{statement}{P3F.36}{Proved}
Let the marked rational elliptic source have either profile
\(I_1^*+I_3+2I_1\) or \(I_1^*+IV+I_1\), and take the
OS50 data \((a,b)=(6,2)\). Put
\(M=\mathcal O(6G-O-L_0)\), with the actual component \(L_0\)
of P3K.21, and place the simple whole-fibre zero at an \(I_1\)
value. At the \(I_1^*\) value the actual primitive quadratic
cap exists for every odd norm numerator \(n\), with its integral
delta phase retained. Leave the other source values ordinary.
The circle orbit of this actual construction satisfies
\[
 \pi_1(Y)=\mathbb Z/|n|,\qquad
 H_j(Y;\mathbb Z)=
 \begin{cases}
 \mathbb Z,&j=0,5,\\
 \mathbb Z/|n|,&j=1,3,\\
 0,&j=2,4.
 \end{cases}
\]
The total six-manifold is a smooth sphere exactly when \(n=1\)
or \(n=-1\). This statement supplies the actual cap and its
integral topology; it assigns no parameter-return class.
\end{statement}

\begin{proof}
\par\medskip\noindent\textbf{The actual root and line correspondence.}
We first give the local construction for each surviving starred
line point \(a=2,6\) and translation \(b\equiv2\pmod4\).
The minimal quadratic root \(S'\to\Delta_z\) has fibre
\(I_2\). Its canonical deck \(D\) combines \(z\mapsto-z\)
with inversion on the Tate fibre. It fixes each of the two central
components and has two fixed smooth points on each. Blow up these
four points by \(\beta:\widehat S\to S'\). The lifted deck
fixes each exceptional curve pointwise; in a chart whose coordinates
before blowing up have signs \((-1,-1)\), the blowup coordinates
\((w/z,z)\) have signs \((1,-1)\), and the quotient coordinates
are \((w/z,z^2)\). The quotient \(\pi:\widehat S\to S\)
is therefore the smooth starred source of P3K.21.

Let \(E_0,E_1\) lie over the leaves \(L_0,L_1\) met by
\(O\) and \(P=aG\). Write \(\widehat O,\widehat P\)
for the strict transforms of the root sections. The actual divisor
pullbacks are
\[
 \pi^*L_0=2E_0,\qquad
 \pi^*(P-O-L_0)=\widehat P-\widehat O-2E_0,
\]
\[
 \beta^*(P'-O')=\widehat P-\widehat O+E_1-E_0.
\]
Consequently, with \(N=\mathcal O(P'-O')\), we have the exact
line comparison
\[
                  \beta^*N=\pi^*M+E_0+E_1.                 \tag{P3F.36.1}
\]
The two root sections meet the same \(I_2\) component, so \(N\)
has degree zero on both components. Translation by \(bG\) swaps
\(E_0,E_1\), hence preserves their sum. Thus (P3F.36.1) transports
the specified suspension Hom line and its lift, including the
vertical correction.

More explicitly, write the Tate period and generator as
\[
 q_T=z^2u(z^2),\qquad Z_G=zv(z),\qquad
                  v(z)v(-z)=-u(z^2).
\]
The degree-zero root point is represented by the nonvanishing unit
\[
 Z_P=v(z)^a/u(z^2)^{a/2},\qquad
 Z_P(-z)=Z_P(z)^{-1},\qquad Z_P(0)=-1.
\]
Its line has Tate transitions
\((Z,\lambda)\sim(q_TZ,Z_P\lambda)\). These unit transitions
extend over the two-component model and realize precisely the
line in (P3F.36.1). They also retain its nonzero order-two central
Picard character.

\par\medskip\noindent\textbf{The actual norm lift and its sign.}
The positive global Hom section gives the contracting map covering
\(t_{-bG}\). Let \(F\) be a line lift on \(N\) of
\(t_G\), provided directly by the Tate unit transitions, and
let \(D\) denote the canonical order-two line lift as well.
Every other lift of this translation differs by a base unit.
Write
\(DFD^{-1}=c(z)F\). Since \(D^2=1\),
\(c(z)c(-z)=1\). Dividing \(F\) by a base unit obtained from
the odd logarithm of \(c(z)/c(0)\) makes
\[
                  DFD^{-1}=\chi F,\qquad \chi\in\{1,-1\}.
\]
The actual nonzero-linearization lift is
\(B_0=\beta_0(z)F^{-b}\). Its descent and even \(b\) imply
\(\beta_0(-z)=\beta_0(z)\). Put \(k=-nb/2\) and choose
analytic square roots of the even nonvanishing base units. Then
\[
 A_0=\chi^{-k/2}\beta_0(z)^{n/2}F^kD
 \quad\hbox{satisfies}\quad
 A_0B_0=B_0A_0,\qquad A_0^2=B_0^n.                         \tag{P3F.36.2}
\]
Indeed \((F^kD)^2=\chi^kF^{2k}\); this character factor
is cancelled by the displayed scalar, not omitted. The integral
choice of its phase is retained in the outer line collar.

For the family parameter \(r\), take
\(\lambda=r^2\), \(B_\lambda=\lambda B_0\), and
\(A_r=r^nA_0\). These are actual maps and
\(A_r^2=B_\lambda^n\). Choose Bezout integers \(r_0,s_0\) with \(r_0n+2s_0=1\). The generator
\(H=A_r^{r_0}B_\lambda^{s_0}\) satisfies
\(A_r=H^n\), \(B_\lambda=H^2\). Its source map is exactly
\[
                         h=t_{-bG/2}D.
\]
Its multiplier contains the positive first power of the small
parameter \(r\), times a bounded line map on a compact smaller
root neighbourhood. It is uniformly contracting for small
\(r\ne0\). Thus no nonzero iterate fixes a nonzero line vector,
and the norm-shell fundamental region proves properness and
compactness as in P3K.15. The quotient is smooth and its line-circle
action is free. The source action exchanges the two central
components.

We fix the outer signs explicitly. Normalize the real line
coordinate so that \(B\) sends \(t\) to \(t-1\), and put
\(\theta=\arg z/(2\pi)\). The actual root action is
\[
 A:(\theta,t)\longmapsto(\theta+1/2,t-n/2).
\]
The original outer angular coordinates are
\[
                       x=2\theta,\qquad y=t+n\theta.
                                                               \tag{P3F.36.3}
\]
The positive base meridian \(\alpha\) changes \(x\) by one
at fixed \(y\), while the positive outward suspension circle
\(\gamma\) changes \(y\) by one. The positive cap generator
\(h\) changes \(t\) by \(1/2\). Hence the free cap map is
\[
                         \alpha\longmapsto-nh,
                  \qquad \gamma\longmapsto2h,               \tag{P3F.36.4}
\]
and its killed meridian is \(2\alpha+n\gamma\).
The exact norm-shell coordinate of P3K.15 straightens the bounded
norm units; the remaining source and delta periods stay in the
integral collar coefficients below.

\par\medskip\noindent\textbf{The source complement and its actual chain maps.}
Now take \(b=2\) and either stated source profile. Let \(P\)
be the complement of a small \(I_1^*\) tube. Its elliptic
fibration over the complementary disc has a section, so van Kampen
presents \(\pi_1(P)\) as a quotient of the elliptic lattice by
the singular-fibre relations. It is abelian. The primitive component
lattice and the unimodular source intersection form give
\(H_1(P)=0\); therefore \(\pi_1(P)=0\). The same sequence
gives \(H_2(P)=\mathbb Z^4\), with no other positive reduced
homology.

Use the actual source basis
\((O,F,C_0,C_1,L_1,R_0,R_1,A_{2,1},A_{2,2},G)\).
Its Gram determinant is \(-1\). The dual columns indexed by
\((O,G,A_{2,1},A_{2,2})\) give the complement pairing
\[
 J=\begin{pmatrix}
 -12&12&8&4\\12&-12&-8&-4\\8&-8&-6&-3\\4&-4&-3&-2
 \end{pmatrix}.
\]
Modulo the starred components, translation by \(2G\) has matrix
\[
 T=\begin{pmatrix}
 -1&-2&0&0\\2&3&0&0\\1&2&-1&1\\0&1&-1&0
 \end{pmatrix},\qquad U=T^{-t}.
                                                               \tag{P3F.36.5}
\]
For example \([2G]=-[O]+2[G]+[A_{2,1}]\) and
\([3G]=-2[O]+3[G]+2[A_{2,1}]+[A_{2,2}]\); translation sends
\(A_{2,1}\) to \(-A_{2,1}-A_{2,2}\) and \(A_{2,2}\)
to \(A_{2,1}\). These are the actual section and component
classes for both profiles. Thus no identification of their complex
sources is needed. The certificate checks \(JU=TJ\).
The actual contracting translation is the inverse, with matrices
\(T^{-1},U^{-1}\).

For clarity we give the exact chain comparison to the auxiliary
positive-translation calculation. The diffeomorphism
\[
 T_{t_{-2G}}P\longrightarrow T_{t_{2G}}P,
                     \qquad[p,t]\longmapsto[p,-t]             \tag{P3F.36.6}
\]
fixes source cycles and reverses the suspension circle. The boundary
complex has groups \(\mathbb Z,\mathbb Z^4,\mathbb Z^4,
\mathbb Z\) in degrees zero through three and sole differential
\(d_2=J\). It is the Poincare--Lefschetz cone for \(P\).
Both absolute and relative split complexes have free homology in
even degrees, so the naturality homotopy has no degree-one term.
Taking the mapping-torus cone therefore gives the actual boundary
and inclusion maps. In degrees one, two, three, (P3F.36.6) has
chain matrices
\[
 \Phi_1=\operatorname{diag}(I_4,-1),\quad
 \Phi_2=\operatorname{diag}(I_4,-T^{-1}),\quad
 \Phi_3=\operatorname{diag}(1,-U^{-1}).
\]
The certificate checks both chain identities before computing
homology. For example its positive-model degree-two cycle columns
are
\[
 R_2=\begin{pmatrix}
 -1&0&1\\0&0&0\\-1&0&0\\-1&0&2\\
 -2&3&-2\\2&-3&0\\2&-2&0\\1&-1&0
 \end{pmatrix};
\]
the actual columns are exactly \(\Phi_2^{-1}R_2\).

Write \((\epsilon,r_2,s_2)\) for these boundary classes,
with orders \((2,0,0)\), and \((\tau,g)\) for the exterior
classes, of orders \((2,0)\). Their actual inclusion and fibre
classes are
\[
 i_2=\begin{pmatrix}1&0&-2\\0&0&1\end{pmatrix},\qquad
 F_{\partial}=\epsilon-2r_2,\qquad F_W=\tau.                  \tag{P3F.36.7}
\]
In particular the fibre is the nonzero order-two exterior class.
The cap retracts onto the mapping torus of the \(I_2\) reflection.
Its \(H_2\) is generated by one component-orbit class \(C\),
and the actual fibre specializes to \(2C\). Thus, for the actual
integral collar coefficient \(\kappa\), its boundary map is
\[
             (\epsilon,r_2,s_2)\longmapsto(0,-1,\kappa)C.
\]
The entire degree-two Mayer--Vietoris map is
\[
 D_2=\begin{pmatrix}
 1&0&-2\\0&0&1\\0&1&-\kappa
 \end{pmatrix}.                                               \tag{P3F.36.8}
\]
The first source and target generators have order two and map
identically; the free block is unimodular. This is an integral
isomorphism for the actual \(\kappa\). Equivalently the explicit
basis change \(s_{2,\mathrm{new}}=s_2+\kappa r_2\) makes its
cap column zero and preserves its exterior column. We do not infer
\(\kappa=0\) from an inversion: the exterior inversion can
change the marked cap by a gamma shear.

\par\medskip\noindent\textbf{The remaining maps and fundamental group.}
The cap central graph has generator \(\eta\), reversed by the
affine reflection. Thus
\[
 \pi_1(\mathrm{cap})=
       \langle h,\eta\mid h\eta h^{-1}=\eta^{-1}\rangle,
\]
\(H_1=\mathbb Zh\oplus\mathbb Z/2\eta\), and
\(H_2=H_3=\mathbb Z\). Let \(u,w\) be the vanishing and
graph periods on the source boundary. In the positive outward
\(\gamma\) convention its affine relation is
\(\alpha\gamma\alpha^{-1}=\gamma w^b\). The actual map,
retaining the graph-path integer \(c\), is
\[
 u\mapsto1,\qquad w\mapsto\eta,\qquad
 \alpha\mapsto h^{-n}\eta^c,\qquad
 \gamma\mapsto h^2\eta^{-b/2}.                               \tag{P3F.36.9}
\]
The commutator identity verifies the graph correction. In degree
one, with the target cap signs absorbed into its chosen generators,
the complete matrix for \(b=2\) is
\[
 D_1=\begin{pmatrix}1&c&-1\\0&-n&2\\0&0&1\end{pmatrix}.
\]
Its first source and target entries have order two. Their map is
the identity and the free determinant is \(-n\).

In degree three use the oriented source boundary and the positively
oriented fibre--gamma trace. The exterior row is \((0,1)\).
The cap generator is the trace of the sum of the two components
through one positive \(h\) turn. The norm powers and
(P3F.36.4) give cap row \((-n,2)\). A graph correction is a
path inside the two-dimensional central source and adds no
three-homology class. Thus
\[
                       D_3=\begin{pmatrix}0&1\\n&-2\end{pmatrix}.
\]
In degree zero the map is \((1,-1)^t\); in degree four the
domain is \(\mathbb Z\) and the target is zero. These specify
every nonzero Mayer--Vietoris map and give the stated integral
homology. All matrices, their chain representatives and the
symbolic collar coefficients are executed in
\path{code/part3/i1star_boundary.py} and its JSON output.

For van Kampen, \(\pi_1(W)=\mathbb Z\gamma\). Gluing kills
\(u,w,\alpha\), hence \(\eta=1\), \(h^n=1\), and
\(\gamma=h^2\). Since \(n\) is odd, the resulting group is
exactly \(\mathbb Z/|n|\).

\par\medskip\noindent\textbf{The actual six-dimensional Hopf filling.}
For \(n=\pm1\), \(Y\) is simply connected and has the
integral homology of the five-sphere, hence is the standard smooth
five-sphere by the smooth classification already used in P3K.22.
The simple whole zero still has the same actual cusp chart and
fixed sphere \(F\cong S^2\), with normal circle weights
\((-1,+1)\). On \(Y\setminus\operatorname{int}\nu F\),
Alexander duality gives the cohomology of \(S^2\). The normal
Hopf meridian has Euler value one, so the principal-circle Gysin
sequence gives the homology of \(S^3\), and its homotopy sequence
gives trivial fundamental group. Gluing \(F\times D^4\) along
\(F\times S^3\), using the actual normal \(SU(2)\) framing,
has maps \((1,-1)^t\) in degree zero, \(-1\) in degree two,
\(1\) in degree three, and zero from \(\mathbb Z\) in degree
five. Thus the resulting \(X\) is simply connected with the
integral homology of \(S^6\). The same six-dimensional smooth
classification as P3K.22 gives the standard smooth sphere.

For general odd \(n\), removing a codimension-three fixed sphere
preserves \(\pi_1(Y)\). Its local Hopf meridian still makes the
circle-boundary map onto \(\pi_1(S^1)\), so the free-circle
piece has the same fundamental group as its orbit. The simply
connected fixed tube and its simply connected boundary do not
change that group. Hence \(\pi_1(X)=\mathbb Z/|n|\), which
excludes every other odd numerator from a sphere. All line phases
and collars were retained throughout; no return invariant follows
from this topological calculation alone.
\end{proof}
